\documentclass[12pt]{article}
\usepackage[utf8]{inputenc}
\usepackage[T1]{fontenc}
\usepackage[margin=1in]{geometry}
\usepackage{setspace}
\usepackage{xcolor}
\usepackage[colorlinks=true, urlcolor=blue!60!black, linkcolor=blue!60!black]{hyperref}
\setlength{\parindent}{2em}
\setlength{\parskip}{1.25em}
\renewcommand{\baselinestretch}{1.0}
\begin{document}
\begin{center}
{\LARGE \textit{United States v. Darby}, 312 U.S. 100 (1941)\par}
\vspace{0.8em}
{\large Case Brief --- Four-Page Edition\par}
\end{center}
\vspace{1em}

\section{Memory Jogger}

Overruling \textit{Hammer v.\ Dagenhart} and dismissing the Tenth Amendment as ``but a truism,'' a unanimous Court upheld the Fair Labor Standards Act of 1938: Congress may prohibit interstate shipment of goods produced under substandard labor conditions and directly regulate intrastate production that substantially affects interstate commerce (\href{https://supreme.justia.com/cases/federal/us/312/100/}{312 U.S. 100}).

\section{Detailed Case Facts}

The Fair Labor Standards Act set a minimum wage of 25 cents an hour and a maximum 44-hour week without overtime pay for employees engaged in commerce or in the production of goods for commerce. Section 15(a)(1) barred interstate shipment of goods produced in violation of those standards; \S~15(a)(2) prohibited employing workers in such production in violation of them; and \S~15(a)(5) required employers to keep wage and hour records. Willful violations were crimes.

Fred W.\ Darby's lumber company in Statesboro, Georgia, milled Georgia timber into lumber, intending to ship much of it out of state, and did so. He was indicted for paying mill workers less than 25 cents an hour, working them more than 44 hours without overtime, shipping the lumber interstate, and failing to keep records.

\section{Procedural History}

Darby demurred, arguing that the Act exceeded the commerce power and violated the Fifth and Tenth Amendments. Relying on \textit{Hammer}, the Southern District of Georgia quashed the indictment, reasoning that manufacturing is not commerce and that intent to ship did not bring local production under federal control. The United States appealed directly, and the Supreme Court reversed on February 3, 1941.

\section{Judicial Votes}

Reversed, \textbf{8--0}.
\begin{itemize}
\item \textbf{Opinion of the Court:} Justice Stone, joined by Chief Justice Hughes and Justices Roberts, Black, Reed, Frankfurter, Douglas, and Murphy.
\item \textbf{Vacancy:} Justice McReynolds retired three days before the decision.
\item \textbf{Concurring or dissenting:} None.
\end{itemize}

\section{Holding}

Under the Commerce and Necessary and Proper Clauses, Congress may prohibit interstate shipment of goods made under substandard labor conditions regardless of its motive; \textit{Hammer v.\ Dagenhart} is overruled. Congress may also directly regulate intrastate production of goods intended for interstate commerce, both as a means of enforcing the shipment ban and because substandard production substantially affects interstate commerce. The Tenth Amendment imposes no independent limit on appropriate means of executing a granted power, and the Act satisfies the Fifth Amendment. The ``truism'' language and the claim that motive is beyond review operated partly as broad statements later cabined in \textit{New York}, \textit{Printz}, and \textit{Lopez}.

\section{Analysis of Opinions}

\textbf{Shipment ban.} Manufacturing is not itself commerce, but shipping goods across state lines is a regulation of commerce. Quoting \textit{Gibbons v.\ Ogden}, Stone held the commerce power complete in itself, including the power to prohibit interstate movement of any goods Congress deems injurious, not only lottery tickets or diseased cattle. Congress's motive, to keep commerce from spreading substandard labor conditions through unfair competition, was a policy judgment beyond judicial control. Praising Justice Holmes's ``powerful and now classic dissent,'' the Court overruled \textit{Hammer}.

\textbf{Direct regulation of production.} Production for commerce includes goods made with the intent or expectation that some will be shipped out of state, since factories rarely segregate inventory. Invoking \textit{McCulloch}, Stone upheld \S~15(a)(2) as a necessary and proper means to make the shipment ban effective and, independently, because substandard wages in intrastate production substantially affect interstate commerce by enabling unfair competition, even by small producers whose aggregate effect is substantial. \textit{Carter v.\ Carter Coal} was limited.

\textbf{Tenth Amendment.} Drawing on the ratification debates and Story's \textit{Commentaries}, Stone read the Amendment as declaratory, not a curb on appropriate means.

\textbf{Records and due process.} Recordkeeping is an appropriate enforcement tool, and under \textit{West Coast Hotel} wage and hour laws satisfy due process. There were no separate opinions.

\section{Examples: Future Applications}

\begin{enumerate}
\item \textbf{Same-Side (Valid): Chemical-Exposure Rules.} Congress bars shipment of electronics made in plants exceeding federal toxin limits and penalizes the employers directly. Both measures fit \textit{Darby}.
\item \textbf{Same-Side (Valid): Records for Small Processors.} Small food plants that sell some output interstate must keep sanitation and wage logs for all employees. Congress may regulate commingled production and impose incidental record duties.
\item \textbf{Opposite-Side (Invalid): Guns Near Schools.} Mere possession of a firearm near a school is non-economic and lacks a substantial economic effect, as \textit{United States v.\ Lopez} held.
\item \textbf{Opposite-Side (Invalid): An Insurance Mandate.} \textit{NFIB v.\ Sebelius} held that the commerce power regulates existing activity and does not compel individuals to enter commerce.
\item \textbf{Fence-Sitter: The Local Workshop.} A three-person Georgia shop sells all its tables locally but buys out-of-state varnish. It does not produce for commerce under \textit{Darby}, but aggregation reasoning later expanded in \textit{Wickard} and \textit{McClung} could reach it.
\end{enumerate}

\section{Critique}

\textbf{Originalist critiques.} Randy Barnett and Richard Epstein argue that ``commerce'' originally meant trade and exchange, distinct from manufacturing and labor relations. By abandoning review of pretextual motives, contrary to Marshall's warning in \textit{McCulloch}, and reducing the Tenth Amendment to a truism, \textit{Darby} effectively gave Congress a general police power and ignored the Amendment's protection of state and popular sovereignty.

\textbf{Progressive and doctrinal defenses.} Bruce Ackerman and others see \textit{Darby} as the culmination of the New Deal transformation, aligning doctrine with an integrated economy in which states face a race to the bottom if they act alone on wages and child labor. Revisionist scholars such as Andrew Coan and David Schwartz argue that Founding-era enumerationism was more contested than critics claim, lending historical plausibility to deference.

\section{Key Quotations}

\begin{quote}
``The amendment states but a truism that all is retained which has not been surrendered. There is nothing in the history of its adoption to suggest that it was more than declaratory of the relationship between the national and state governments as it had been established by the Constitution before the amendment, or that its purpose was other than to allay fears that the new national government might seek to exercise powers not granted, and that the states might not be able to exercise fully their reserved powers.''
\end{quote}
\noindent --- Justice Stone, Opinion of the Court, 312 U.S. at 124.

\begin{quote}
``The power of Congress over interstate commerce is not confined to the regulation of commerce among the states. It extends to those activities intrastate which so affect interstate commerce or the exercise of the power of Congress over it as to make regulation of them appropriate means to the attainment of a legitimate end, the exercise of the granted power of Congress to regulate interstate commerce.''
\end{quote}
\noindent --- Justice Stone, Opinion of the Court, 312 U.S. at 118.

\begin{quote}
``Whatever their motive and purpose, regulations of commerce which do not infringe some constitutional prohibition are within the plenary power conferred on Congress by the Commerce Clause.''
\end{quote}
\noindent --- Justice Stone, Opinion of the Court, 312 U.S. at 115.

\end{document}
